% GENERATED FILE. Edit canonical lessons, then run npm run generate:books.
% canonical-source-hash: fnv1a64:0e72d3f0ec18bbd4
% canonical-lessons: TA-C100-children, TA-C100-boy, TA-C100-husband, TA-C100-wife, TA-C100-relative

\chapter{Children, Husbands, Wives}
\label{ch:ta-children-husbands-wives}
\hlchaptermodality{100}

\begin{chapteropening}
\textbf{By the end of this chapter:} \emph{I can talk about children, a boy, a husband or wife, and a relative.}

Complete the last lesson of chapter 100: I can talk about children, a boy, a husband or wife, and a relative.
\end{chapteropening}

\section[\texorpdfstring{\ta{குழந்தைகள்} (kuḻantaikaḷ)}{குழந்தைகள் (kuḻantaikaḷ)}]{\ta{குழந்தைகள்} (kuḻantaikaḷ) — children}

\label{lesson:TA-C100-children}



\begin{quote}
Before the new one: say the Tamil for please forgive me, then the Tamil for don't worry.
\end{quote}

\subsection*{You'll want to know: \ta{குழந்தைகள்}}
\textbf{\ta{குழந்தைகள்}} (\emph{kuḻantaikaḷ}) — \textquotedblleft{}children\textquotedblright{}.

\textbf{-\ta{கள்}} is the plural. \textbf{\ta{உங்களுக்குக்} \ta{குழந்தைகள்} \ta{இருக்கிறார்களா}?} — \textquotedblleft{}do you have children?\textquotedblright{}.

\textbf{In the family, and next door}

\noindent
\begin{tabularx}{\linewidth}{@{}>{\raggedright\arraybackslash}X>{\raggedright\arraybackslash}X@{}}
\toprule
\textbf{} & \textbf{word} \\
\midrule
Telugu & \te{పిల్లలు} (\emph{pillalu}) \\
Hindi (neighbour) & \dv{बच्चे} (\emph{bacce}) \\
English & children \\
\bottomrule
\end{tabularx}

\begin{grammarlens}[title={What you've built}]
The question that follows a name.
\end{grammarlens}

\subsection*{Your turn}
\begin{itemize}
\raggedright
  \item \emph{Say it:} \emph{kuḻantaikaḷ}
  \item \emph{Say it:} \emph{kuḻantaikaḷ}, once more
  \item \emph{Say it:} ask about children
  \item \emph{Recall:} say the Tamil for please forgive me, then the Tamil for don't worry, then say \emph{kuḻantaikaḷ} again
\end{itemize}

\begin{tcolorbox}[breakable,colback=teal!4,colframe=teal!35!black,title={Before you move on}]
What does \ta{குழந்தைகள்} mean? (\textquotedblleft{}Children\textquotedblright{}.) Say it once more.
\end{tcolorbox}
\section[\texorpdfstring{\ta{பையன்} (paiyaṉ)}{பையன் (paiyaṉ)}]{\ta{பையன்} (paiyaṉ) — a boy; a son}

\label{lesson:TA-C100-boy}



\begin{quote}
Before the new one: say the Tamil for don't worry, then the Tamil for children.
\end{quote}

\subsection*{You'll want to know: \ta{பையன்}}
\textbf{\ta{பையன்}} (\emph{paiyaṉ}) — \textquotedblleft{}a boy\textquotedblright{}, and in a family, \textquotedblleft{}a son\textquotedblright{}.

\textbf{\ta{என்} \ta{பையன்}} — \textquotedblleft{}my boy\textquotedblright{}, \textquotedblleft{}my son\textquotedblright{}.

\textbf{In the family, and next door}

\noindent
\begin{tabularx}{\linewidth}{@{}>{\raggedright\arraybackslash}X>{\raggedright\arraybackslash}X@{}}
\toprule
\textbf{} & \textbf{word} \\
\midrule
Telugu & \te{అబ్బాయి} (\emph{abbāyi}) \\
Hindi (neighbour) & \dv{लड़का} (\emph{laṛkā}) \\
English & a boy \\
\bottomrule
\end{tabularx}

\begin{grammarlens}[title={What you've built}]
A boy.
\end{grammarlens}

\subsection*{Your turn}
\begin{itemize}
\raggedright
  \item \emph{Say it:} \emph{paiyaṉ}
  \item \emph{Say it:} \emph{paiyaṉ}, once more
  \item \emph{Say it:} say \emph{eṉ paiyaṉ}
  \item \emph{Recall:} say the Tamil for don't worry, then the Tamil for children, then say \emph{paiyaṉ} again
\end{itemize}

\begin{tcolorbox}[breakable,colback=teal!4,colframe=teal!35!black,title={Before you move on}]
What does \ta{பையன்} mean? (\textquotedblleft{}A boy; a son\textquotedblright{}.) Say it once more.
\end{tcolorbox}
\section[\texorpdfstring{\ta{கணவர்} (kaṇavar)}{கணவர் (kaṇavar)}]{\ta{கணவர்} (kaṇavar) — a husband}

\label{lesson:TA-C100-husband}



\begin{quote}
Before the new one: say the Tamil for children, then the Tamil for a boy.
\end{quote}

\subsection*{You'll want to know: \ta{கணவர்}}
\textbf{\ta{கணவர்}} (\emph{kaṇavar}) — \textquotedblleft{}a husband\textquotedblright{}, in its respectful form.

\textbf{\ta{என்} \ta{கணவர்}} — \textquotedblleft{}my husband\textquotedblright{}.

\textbf{In the family, and next door}

\noindent
\begin{tabularx}{\linewidth}{@{}>{\raggedright\arraybackslash}X>{\raggedright\arraybackslash}X@{}}
\toprule
\textbf{} & \textbf{word} \\
\midrule
Kannada & \ka{ಗಂಡ} (\emph{gaṇḍa}) \\
Telugu & \te{భర్త} (\emph{bharta}) \\
Malayalam & \ml{ഭർത്താവ്} (\emph{bharttāvŭ}) \\
Hindi (neighbour) & \dv{पति} (\emph{pati}) \\
English & a husband \\
\bottomrule
\end{tabularx}

\begin{grammarlens}[title={What you've built}]
A husband.
\end{grammarlens}

\subsection*{Your turn}
\begin{itemize}
\raggedright
  \item \emph{Say it:} \emph{kaṇavar}
  \item \emph{Say it:} \emph{kaṇavar}, once more
  \item \emph{Say it:} say \emph{eṉ kaṇavar}
  \item \emph{Recall:} say the Tamil for children, then the Tamil for a boy, then say \emph{kaṇavar} again
\end{itemize}

\begin{tcolorbox}[breakable,colback=teal!4,colframe=teal!35!black,title={Before you move on}]
What does \ta{கணவர்} mean? (\textquotedblleft{}A husband\textquotedblright{}.) Say it once more.
\end{tcolorbox}
\section[\texorpdfstring{\ta{மனைவி} (maṉaivi)}{மனைவி (maṉaivi)}]{\ta{மனைவி} (maṉaivi) — a wife}

\label{lesson:TA-C100-wife}



\begin{quote}
Before the new one: say the Tamil for a boy, then the Tamil for a husband.
\end{quote}

\subsection*{You'll want to know: \ta{மனைவி}}
\textbf{\ta{மனைவி}} (\emph{maṉaivi}) — \textquotedblleft{}a wife\textquotedblright{}.

It shares its start with \textbf{\ta{மனை}}, \textquotedblleft{}a house\textquotedblright{}, from the hospital word: the wife is named from the home.

\textbf{In the family, and next door}

\noindent
\begin{tabularx}{\linewidth}{@{}>{\raggedright\arraybackslash}X>{\raggedright\arraybackslash}X@{}}
\toprule
\textbf{} & \textbf{word} \\
\midrule
Kannada & \ka{ಹೆಂಡತಿ} (\emph{heṇḍati}) \\
Telugu & \te{భార్య} (\emph{bhārya}) \\
Malayalam & \ml{ഭാര്യ} (\emph{bhārya}) \\
Hindi (neighbour) & \dv{पत्नी} (\emph{patnī}) \\
English & a wife \\
\bottomrule
\end{tabularx}

\begin{grammarlens}[title={What you've built}]
Husband and wife.
\end{grammarlens}

\subsection*{Your turn}
\begin{itemize}
\raggedright
  \item \emph{Say it:} \emph{maṉaivi}
  \item \emph{Say it:} \emph{maṉaivi}, once more
  \item \emph{Say it:} say \emph{eṉ maṉaivi}
  \item \emph{Recall:} say the Tamil for a boy, then the Tamil for a husband, then say \emph{maṉaivi} again
\end{itemize}

\begin{tcolorbox}[breakable,colback=teal!4,colframe=teal!35!black,title={Before you move on}]
What does \ta{மனைவி} mean? (\textquotedblleft{}A wife\textquotedblright{}.) Say it once more.
\end{tcolorbox}
\section[\texorpdfstring{\ta{உறவினர்} (uṟaviṉar)}{உறவினர் (uṟaviṉar)}]{\ta{உறவினர்} (uṟaviṉar) — a relative}

\label{lesson:TA-C100-relative}



\begin{quote}
Before the new one: say the Tamil for a husband, then the Tamil for a wife.
\end{quote}

\subsection*{You'll want to know: \ta{உறவினர்}}
\textbf{\ta{உறவினர்}} (\emph{uṟaviṉar}) — \textquotedblleft{}a relative\textquotedblright{}.

\textbf{\ta{உறவு}} is \textquotedblleft{}a tie, a relationship\textquotedblright{}; a relative is one who has one with you.

\textbf{In the family, and next door}

\noindent
\begin{tabularx}{\linewidth}{@{}>{\raggedright\arraybackslash}X>{\raggedright\arraybackslash}X@{}}
\toprule
\textbf{} & \textbf{word} \\
\midrule
Kannada & \ka{ಸಂಬಂಧಿಕ} (\emph{sambandhika}) \\
Telugu & \te{బంధువు} (\emph{bandhuvu}) \\
Malayalam & \ml{ബന്ധു} (\emph{bandhu}) \\
Hindi (neighbour) & \dv{रिश्तेदार} (\emph{riśtedār}) \\
English & a relative \\
\bottomrule
\end{tabularx}

\begin{grammarlens}[title={What you've built}]
Children, boy, husband, wife, relative.
\end{grammarlens}

\subsection*{Your turn}
\begin{itemize}
\raggedright
  \item \emph{Say it:} \emph{uṟaviṉar}
  \item \emph{Say it:} \emph{uṟaviṉar}, once more
  \item \emph{Say it:} say \emph{avar eṉ uṟaviṉar}
  \item \emph{Recall:} say the Tamil for a husband, then the Tamil for a wife, then say \emph{uṟaviṉar} again
\end{itemize}

\begin{tcolorbox}[breakable,colback=teal!4,colframe=teal!35!black,title={Before you move on}]
What does \ta{உறவினர்} mean? (\textquotedblleft{}A relative\textquotedblright{}.) Say it once more.
\end{tcolorbox}
